\section{Comparison with CDP20.}

In this section we compare our construction with the non-existence results of Campana, Demailly, and Peternell [CDP20] (arXiv:2005.13523).

\subsection{Analysis of the CDP20 deformation obstruction}

In [CDP20, Proposition~2.4], the authors consider a smooth complex $3$-manifold $M$ fibered over $\dbP^1$ with abelian surface fibers, and argue that under certain hypotheses, deformations of $M$ are obstructed.
Specifically, Proposition~2.4 of [CDP20] relies on:
\begin{quote}
\textbf{Hypothesis (1) of [CDP20, Prop.~2.4]:}
For every line bundle $L \in \Pic(M)$, the second direct image sheaf vanishes at the critical point $p_0$:
\[
(R^2 f_*(TM \otimes L))_{p_0} = 0.
\]
\end{quote}

\begin{theorem}[Failure of CDP20 Hypotheses on $X$]\label{thm:cdp-failure}
The complex $3$-manifold $X$ constructed in Section~6 violates Hypothesis~(1) of [CDP20, Prop.~2.4] for \textbf{every} line bundle $L \in \Pic(X)$.
Specifically:
\begin{enumerate}[label=(\roman*)]
    \item The central fiber $W_0 = f^{-1}(p_0)$ is a \textbf{non-normal} complex analytic surface, whose singular locus $D = \mathrm{Sing}(W_0)$ has complex codimension $1$ in $W_0$.
    \item By Serre--Grothendieck duality on the Cohen--Macaulay space $W_0$:
    \[
    (R^2 f_*(TX \otimes L))_{p_0} \cong H^2(W_0, TX|_{W_0} \otimes L) \cong \operatorname{Hom}_{W_0}(TX|_{W_0} \otimes L, \, \omega_{W_0})^\vee.
    \]
    \item The differential of the fibration $df \colon TX|_{W_0} \to \cO_{W_0}$ induces a non-zero section of the conductor sheaf:
    \[
    s = df \otimes e \in H^0(W_0, \, \mathscr{C}_{W_0} \otimes (TX|_{W_0})^\vee).
    \]
    Since the singular locus $D$ has codimension $1$, the Riemann extension theorem fails, and $s$ does not vanish on $D$.
    \item Consequently, $(R^2 f_*(TX \otimes L))_{p_0} \ne 0$ for all $L$, so [CDP20, Prop.~2.4] is \textbf{vacuous} for $X$.
\end{enumerate}
\end{theorem}

\begin{theorem}[Singular Fiber Count Reconciliation]\label{thm:cdp-count}
The singular fiber count formula in [CDP20, Lemma~4.2] asserted $r \le 2$ under the assumption of normal fibers.
Taking into account the parabolic monodromy coinvariants of $T_0$ on the non-normal central fiber $W_0$, the corrected formula gives:
\[
r = s - 1 + t' = 3 - 1 + 1 = 3,
\]
matching exactly the three singular fibers $W_0, 3S_1, 4S_2$ of $f \colon X \to \dbP^1$.
\end{theorem}
